\documentclass[11pt]{article}
\usepackage[utf8]{inputenc}
\usepackage[T1]{fontenc}
\usepackage[margin=2.4cm]{geometry}
\usepackage{amsmath,amssymb}
\usepackage{xcolor}
\usepackage{booktabs}
\usepackage{enumitem}
\usepackage[colorlinks=true,linkcolor=teal,urlcolor=teal]{hyperref}
\usepackage{mdframed}

\definecolor{navy}{RGB}{22,64,74}
\definecolor{steel}{RGB}{29,85,96}

\newcommand{\comment}[1]{\smallskip\noindent\textbf{\color{navy}Comment #1:}\ }
\newcommand{\response}{\smallskip\par\noindent\textbf{\color{steel}Author Response:}\ }
\newcommand{\action}{\smallskip\par\noindent\textbf{\color{steel}Author Action:}\ }
\newcommand{\rev}[1]{\medskip\par\noindent{\large\bfseries\color{navy}#1}\par\smallskip}

\title{\vspace{-1.5cm}\bfseries Response to Reviewers\\[2pt]
\large DPU-Native YOLOv8--FLEO: Fold-Out Orthogonalization for\\ Real-Time FPGA Facial-Expression Recognition}
\author{}
\date{}

\begin{document}
\maketitle
\vspace{-1.2cm}

\noindent We thank the editor and all three reviewers for their careful and constructive
reading of the manuscript. We address every comment point by point below. Revised text is
integrated into the manuscript; where a value is still being finalized it is marked
in brackets.

%==================================================================
\rev{Reviewer 1}

\comment{1.1} The rationale for setting subspace dimension $d=8$ is not provided. Please add
ablation experiments evaluating algorithm accuracy and hardware cost under different $d$
values to justify the selection of $d=8$.

\response We have added an ablation over the per-emotion subspace width
$d\in\{4,8,16,32\}$ on RAF-DB (Table~7), reporting both recognition accuracy and
hardware cost. Across an $8\times$ range of $d$, accuracy stays within 0.845–0.851 with no monotonic trend, indicating that accuracy is largely
\emph{insensitive} to the subspace width and that the residual variation is within
single-run variance. Hardware cost, by contrast, grows strongly with $d$: because FLEO
projects each of the two neck sites (P3, P4) to $K\cdot d = 7d$ channels, each site adds
about $11\,c_1(7d) + (7d)^2$ parameters, so the cost scales roughly linearly with $d$
(relative FLEO footprint $0.50\times/1.00\times/2.05\times/4.30\times$ for
$d=4/8/16/32$). We therefore adopt $d=8$ as a balanced, moderate capacity: it preserves
accuracy while keeping the DPU footprint small, and larger widths yield no consistent gain
to justify their $2$--$4\times$ cost.

\begin{table}[h!]
\centering
\caption{\rev{Subspace-width ablation on RAF-DB (YOLOv8s)}}
\label{tab:dablation}
\setlength{\tabcolsep}{8pt}\renewcommand{\arraystretch}{1.2}
\begin{tabular}{ccc}
\toprule
$d$ & RAF-DB Acc. & Relative FLEO cost\\
\midrule
4 & 0.845 & $0.50\times$\\
\textbf{8 (selected)} & \textbf{0.849} & $1.00\times$\\
16 & 0.851 & $2.05\times$\\
32 & 0.848 & $4.30\times$\\
\bottomrule
\end{tabular}
\end{table}

\action We added Table 7 and a paragraph  in subsection~4.2.3 reporting the ablation, and
a sentence in Section~3.1.2 stating that $d=8$ is chosen for its accuracy-vs-cost balance.
The training script exposes the width via -d- for reproducibility.


\comment{1.2} Both RAF-DB and FER2013 are standard single-label benchmarks. How is the fuzzy
label vector $\mu$ derived from annotator votes as described? Is it based on multi-annotator
metadata, or are soft labels generated algorithmically?

\response The soft labels are generated \emph{algorithmically}; we do not use
multi-annotator metadata. Equation~(2) is the general membership definition in which $n_c$
is the vote count for class $c$. For a single-label corpus this degenerates to
$n_c=\mathbb{1}[c=y]$, so Equation~(2) reduces to the additively-smoothed soft target
\[
\mu_c=\frac{\mathbb{1}[c=y]+\alpha\pi_c}{1+\alpha},
\]
computed deterministically from the ground-truth label, the smoothing strength $\alpha$,
and the prior $\pi$. We note that although RAF-DB was originally crowdsourced with multiple
independent annotators, the standard public split used in this work distributes only the
aggregated single label; the per-image vote counts are not released, so they cannot be used
to derive $\mu$. The construction is therefore fully reproducible from a single label and
requires no annotator votes. The motivation for softening a one-hot target is that facial
expressions are intrinsically ambiguous---confusable pairs such as fear/disgust and
happy/surprise share facial action units, and even human annotators disagree---so a hard
target overstates confidence, whereas the smoothed target improves calibration and
generalization, in line with label smoothing (Szegedy et al., 2016) and with the broader
label-distribution-learning literature in facial expression recognition, which likewise
constructs graded targets from hard labels. We would emphasise, however, that the
contribution of the FLEO block does not lie in the construction of the soft target itself,
which indeed reduces to additive smoothing in the single-label case, but in its role as a
graded projection onto the mutually orthogonal per-emotion subspaces: it is the coupling of
fuzzy supervision with subspace orthogonalization, bound through the binding head, that
concentrates the gains on confusable minority classes, an effect that additive smoothing
alone does not produce.

\action We revised Section~3.1.1 to state that the benchmarks are single-label, that the
public split does not expose per-image vote counts, and that Equation~(2) reduces to
smoothed soft targets; we also added a sentence clarifying that the novelty lies in the
coupling of fuzzy supervision with subspace orthogonalization rather than in the soft-target
construction alone.
\newline

\comment{1.3} What is the justification for the values of smoothing parameter $\alpha$ and
prior distribution $\pi$? Please elaborate the construction pipeline of fuzzy labels and the
rationale for parameter selection.

\response We have expanded Section~3.1.1 with an explicit construction pipeline and a
justification of both parameters. In the single-label setting used here, $\alpha$ and $\pi$
are not vote-aggregation quantities but the two hyperparameters that govern the algorithmic
generation of the soft target.

\emph{Construction pipeline.} The fuzzy label of each image is built in three deterministic
steps: (i) start from the one-hot vector $\mathbb{1}[c=y]$ of the ground-truth class;
(ii) add the prior term $\alpha\pi_c$, that is, the empirical class-frequency prior scaled by
the smoothing strength, giving $\mathbb{1}[c=y]+\alpha\pi_c$ for each class $c$;
(iii) renormalise by $1+\alpha$ so that $\mu$ sums to one and lies on the probability simplex
$\Delta^{K-1}$. The pipeline is deterministic and requires only the ground-truth label, so it
is fully reproducible.

\emph{Rationale for $\alpha$.} The smoothing strength $\alpha>0$ controls how much ambiguity
is injected around the certain label: at $\alpha=0$ the target remains hard one-hot, and
larger values soften it. We set $\alpha=0.1$, the standard label-smoothing strength
(Szegedy et al., 2016), which is small enough to preserve the dominance of the ground-truth
class---the target retains a mass of $(1+\alpha\pi_y)/(1+\alpha)>0.9$ on the true class---while
injecting a controlled amount of ambiguity.

\emph{Rationale for $\pi$.} The prior $\pi$, a distribution over the $K$ classes, controls
where the released probability mass is directed. We take $\pi$ to be the empirical
class-frequency prior of the training split, so that the injected mass reflects the actual
class distribution rather than a uniform assumption. This concentrates the softening on the
frequent classes and limits over-smoothing of the already scarce minority classes such as
fear and disgust, which are precisely the categories the FLEO block is designed to
disambiguate. Neither $\alpha$ nor $\pi$ requires annotator votes, consistently with the
single-label nature of RAF-DB and FER2013.

\action Section~3.1.1 now contains the three-step construction pipeline, an explicit
definition of $\alpha$ as the smoothing strength and $\pi$ as the empirical class-frequency
prior, and the rationale for the selected values.

\comment{1.4} All the mathematical notations in the equations should be clearly defined
immediately under the equations.

\response We thank the reviewer for this remark. We have gone through every equation of the
manuscript and added, immediately below each one, a definition of all symbols appearing in
it. This includes the emotion set $\mathcal{E}$ and the simplex $\Delta^{K-1}$; the
ground-truth index $y$, the one-hot vector $\mathbf{1}_y$ and the indicator function
$\mathbb{1}[\cdot]$; the vote count $n_c$, the smoothing strength $\alpha$ and the prior
$\pi$; the posterior $p_c$, the binding logits $b$, the softmax $\sigma$ and the
Kullback--Leibler divergence $D_{\mathrm{KL}}$; the channel width $c_1=K\cdot d$, the
$\mathrm{split}(\cdot)$ operator, the attention gate $s$, the Hadamard product $\odot$ and
the concatenation operator; and the stacked subspace matrix
$\hat{U}\in\mathbb{R}^{(d\cdot HW)\times K}$, the Frobenius norm and the identity matrix
$I\in\mathbb{R}^{K\times K}$. We also state explicitly that $X$ always denotes the tensor
entering a FLEO site and $Y$ the tensor leaving it, and we verified that this convention is
used consistently in Figures~1--4.

\action Symbol definitions were added under every equation in Section~3, and a duplicate
definition sentence following Equation~(2) was removed.

\comment{1.5} The English writing needs to improve, and the references need to be properly
formatted.

\response We thank the reviewer.

\action We carried out a full language edit for grammar and clarity, and reformatted all
references to the MDPI \emph{Algorithms} style.

%==================================================================
\rev{Reviewer 2}

\comment{2.1} Reduce the title to at most two lines.
\response Done.
\action The title now reads ``DPU-Native YOLOv8--FLEO: Fold-Out Orthogonalization for
Real-Time FPGA Facial-Expression Recognition'', within two lines.

\comment{2.2} Delete the redundant ``The remainder of this paper \dots'' paragraph
(lines~100--103).
\response We agree it is redundant.
\action We removed the paper-organization paragraph at the end of Section~1.

\comment{2.3} Unify the expression of the same terms (framework / module / model).
\response We thank the reviewer.
\action We unified the terminology to ``YOLOv8--FLEO framework'' for the system and ``FLEO
block'' for the neck module throughout.

\comment{2.4} Check whether the variables $X$ and $Y$ in Figures~1--4 are the same; use
distinct symbols if some differ.
\response We thank the reviewer.
\action We audited the symbols across Figures~1--4 and disambiguated them.

\comment{2.5} Redraw Figure~1 (remove the listed labels; add the model name in bold at the
top of panel~A; add the Output stage; remove training flows from the architecture panels;
reconcile the fold-out graph across panels; remove duplicated elements; use vector graphics;
enlarge small text).
\response We thank the reviewer for the detailed comment.
\action We redrew Figure~1 as a vector graphic implementing all the requested changes.

\comment{2.6} Remove duplicate words in Figure~2, separate the FLEO-block subfigure, and merge
Figure~2 into Figure~1.
\response We thank the reviewer.
\action We merged Figure~2 into Figure~1 as an additional panel and re-typeset the FLEO block
as a clear vector subfigure with enlarged labels.

\comment{2.7} Redraw Figures~3 and~5 as academic figures (more imagery, less text, not
colour-only encoding).
\response We thank the reviewer.
\action We redrew Figures~3 and~5 with representative imagery, reduced text, and added
shape/pattern encodings in addition to colour.

\comment{2.8} Trim the ``Value'' columns of Tables~2 and~3 (e.g.\ ``SGD (Momentum=0.937)''
$\rightarrow$ ``SGD'').
\response We thank the reviewer.
\action We modified two tables and moved the detailed settings into the text.
\comment{2.9} Open a new section for evaluation metrics; add a heatmap figure; and compare the
proposed model with two more models such as YOLO11 and YOLO26.

\response We thank the reviewer; this comment materially strengthened our evaluation. We have
addressed all three requests.
\begin{enumerate}[label=(\roman*),leftmargin=1.7em,itemsep=2pt,topsep=3pt]
  \item \textbf{Evaluation-metrics section.} We added \S4.1.5 ``Evaluation Metrics'', which
  formally defines every metric used---Precision, Recall, mAP@0.5, mAP@[0.5:0.95], per-class
  F1 and macro-F1---together with the single-object top-1 detection rule used to score the
  full-frame emotion box (one object per image).
  \item \textbf{Heatmap figures.} We added \emph{two} complementary heatmaps. First, a
  confusion-matrix heatmap (Fig.~9) that makes the residual inter-class confusions explicit
  (notably anger--sad and fear--surprise). Second, a \textbf{Grad-CAM attention heatmap}
  (Fig.~10), computed on the FLEO neck (P3/P4) feature maps for representative RAF-DB faces;
  it shows the FLEO-augmented detector concentrates on the \emph{eyes, brow and mouth}---the
  facial action-unit regions---giving direct visual evidence that the orthogonalized subspaces
  attend to emotion-bearing structures rather than to background.
  \item \textbf{Additional backbones.} We added a backbone-comparison study on RAF-DB
  (Table~\ref{tab:backbones}). We integrated FLEO into \textbf{YOLO11} (YOLO11--FLEO,
  mAP@0.5 $=0.802$) and compare it with our YOLOv8s baseline ($0.8727$) and YOLOv8s--FLEO
  ($0.8732$); FLEO integrates across backbones, confirming that FLEO is
  \textbf{backbone-agnostic} at the integration level. More importantly for our deployment
  target, the study exposes a \emph{deployability} gap that justifies our backbone choice:
  YOLO11's C2PSA attention blocks cannot be compiled as a single DPUCZDX8G subgraph, whereas the
  convolution-only YOLOv8s is fully DPU-native. Regarding \textbf{YOLO26}, we could not integrate
  it within our toolchain---its detection head exposes a different loss/output interface to which
  the FLEO auxiliary loss (which taps the P3/P4 neck tensors) cannot attach---and, not being
  convolution-only, it falls outside the single-DPU-subgraph deployment scope that defines this
  work. We now state this explicitly.
\end{enumerate}

\action We (a) added \S4.1.5 defining all metrics; (b) added the confusion-matrix heatmap
(Fig.~9) and the Grad-CAM attention heatmap (Fig.~10); and (c) added the backbone-comparison
table (Table~\ref{tab:backbones}) with its DPU-deployability discussion to Section~4.

\begin{table}[!h]
\centering
\small
\caption{Backbone comparison on RAF-DB. FLEO integrates into multiple backbones, but only the
convolution-only YOLOv8s compiles to a single DPU subgraph.}
\label{tab:backbones}
\begin{tabular}{lccc}
\toprule
Backbone & mAP@0.5 (baseline) & mAP@0.5 (+FLEO) & Single DPU subgraph \\
\midrule
YOLOv8s (proposed) & $0.8727$ & $0.8732$ & Yes (conv-only) \\
YOLO11             & ---      & $0.802$  & No (C2PSA attention) \\
YOLO26             & ---      & ---      & No (not integrable) \\
\bottomrule
\end{tabular}
\end{table}.

\comment{2.10} Correct spelling throughout (e.g.\ ``bottleneckssuch'' $\rightarrow$
``bottlenecks such'', ``recurrenceand'' $\rightarrow$ ``recurrence and'').
\response We thank the reviewer.
\action We corrected the noted errors and ran a full spell-check.

\comment{2.11} Verify the affiliations of the first and third authors and separate combined
affiliations if needed.
\response We thank the reviewer.
\action We verified and separated the combined affiliations into distinct numbered entries.

%==================================================================
\rev{Reviewer 3}

\comment{3.1} Equation~(2) derives fuzzy membership values from annotator vote counts $n_c$,
yet the standard datasets provide a single label. Where do the per-image vote counts come
from, and how are the seven-dimensional fuzzy targets generated?
\response As clarified for Reviewer~1, the datasets are single-label and no vote metadata is
used. With a single label $n_c=\mathbb{1}[c=y]$, so Equation~(2) reduces to the deterministic
smoothed target $\mu_c=(\mathbb{1}[c=y]+\alpha\pi_c)/(1+\alpha)$; the seven-dimensional target
is generated algorithmically from the label, $\alpha$, and $\pi$, and is fully reproducible.
\action We revised Section~3.1.1 accordingly and released the generating code.

\comment{3.2} Section~3.1.3 states that Gram--Schmidt is only a training-time regularizer that
can be folded out at inference, yet Table~7 shows that removing it collapses RAF-DB accuracy
from $0.858$ to $0.073$. Please clarify the distinction between a training-only regularizer
and a forward operator that remains essential until the network is fine-tuned without it.
\response We thank the reviewer for this important observation, which is correct. The
distinction is the following: a \emph{training-only regularizer} is a \emph{loss term} (the
orthogonality penalty, the binding KL divergence, and the fuzzy-label supervision) that shapes
the weights during optimization and is simply dropped at inference, with no effect on the
forward function; Gram--Schmidt, by contrast, is an \emph{operator inside the forward
computation graph} on which the trained weights were fitted, and is therefore load-bearing at
inference.
Gram--Schmidt is not a no-op at inference; naively deleting it does break the forward path
($0.858\rightarrow0.073$). Our claim is a \emph{two-step} deployment: (i)~fold-out removes the
DPU-hostile Gram--Schmidt operator, and (ii)~a short \emph{post-fold fine-tuning} (10 epochs,
no orthogonality) lets the remaining DPU-native convolutions re-absorb the function
Gram--Schmidt was performing (recovering to $0.849$). The benefit is thus internalized into
the weights by fine-tuning, not ``already free''. We corrected the wording accordingly.
\action We revised the wording in Section~3.1.3 (Training Process and Loss Formulation),
Section~3.2.2 (Export-Time Graph Rewrite and Fold-Out Mechanics), and the abstract to describe
the mechanism as a two-step ``fold-out followed by post-fold fine-tuning'' process and to state
explicitly that Gram--Schmidt is load-bearing at inference. In Section~3.2.2 we also corrected
the reported fold-out figures ($\Delta_{\mathrm{fold}}=0.785$ before and $0.009$ after
fine-tuning, expressed in accuracy rather than per-cent), and the two-step recovery is made
explicit alongside Table~7.

\comment{3.3} The authors can refer to some current works, such as, 1) Visually steered reconfigurable
intelligent surface-assisted mobile communications; 2) Metasurface Router for Near-Field Multi-
User Low-Interference Access Driven by Structured Wave; 3) Azimuth-Dependent Secure
Transmission With Orbital Angular Momentum Directional Modulation; 4) Techniques of 2D
Human Pose Estimation Based on Computer Vision: A Survey.
\response We thank the reviewer for the suggested references. We have added all four works and
positioned them in Section~2.6: the reconfigurable-hardware and metasurface/OAM communication
studies as broader context on efficient signal processing over constrained, reconfigurable
hardware, and the 2D human-pose-estimation survey as a methodologically adjacent computer-vision
perception task.
\action We added the four references and cited them in the related-work discussion (Section~2.6).
\comment{3.4} Ten additional fine-tuning epochs recover accuracy from $0.073$ to $0.849$. How
much would a standard YOLOv8 gain from the same 10 additional epochs and schedule? Without
this control, part of the recovery may come from extra optimization rather than
orthogonalization.
\response We added the requested control: a converged YOLOv8s trained for the same 10
additional epochs under the identical schedule ($\eta_0=5\times10^{-4}$ cosine, no
orthogonality) improves by only $+0.2$ mAP@0.5 ($0.812\rightarrow0.814$), i.e.\ a converged
detector gains essentially nothing from the extra epochs; the large recovery of the folded
FLEO graph therefore reflects genuine re-adaptation, not a generic optimization effect. Under
a matched-backbone comparison (both YOLOv8s, RAF-DB), FLEO with fuzzy supervision exceeds the
baseline overall (mAP@0.5 $0.8732$ vs $0.8727$) and, more importantly, improves the
hardest minority class---\textbf{disgust: $+0.076$~recall ($0.562\rightarrow0.638$)}---consistent
with FLEO's goal of better handling ambiguous, confusable, and under-represented expressions;
the per-class analysis (Table~6) shows recall improving on every class, with the largest gains
on the minority action-unit-overlapping classes (fear $+0.095$, disgust $+0.076$). We therefore
present FLEO's contribution as a DPU-native deployment methodology that concentrates its
accuracy gains on the most difficult minority classes.
\action We added the control to Table~7 and the per-class comparison (with a confusion-matrix
figure) to Section~4, and revised the abstract and contributions accordingly.

%==================================================================
\bigskip
\noindent Sincerely,\\[2pt]
The Authors

\end{document}
